\documentclass[runningheads]{llncs}

\usepackage[T1]{fontenc}
\usepackage{newtxtext}
\usepackage[varvw]{newtxmath}
\usepackage{graphicx}
\usepackage{booktabs}
\usepackage{multirow}
\usepackage{url}

\title{ERRORFlow: Training-Free Selective Residual Correction for Fact Verification}
\titlerunning{ERRORFlow}

% TODO: replace author and affiliation placeholders before submission.
\author{Anonymous Author(s)}
\authorrunning{Anonymous Author(s)}
\institute{Anonymous Institution}

\begin{document}
\maketitle

\begin{abstract}
Large language model fact verifiers can benefit from a second inference pass,
but applying correction prompts uniformly increases compute and may regress
already-correct predictions. We introduce ERRORFlow, a training-free residual
correction framework that stores source-domain action outcomes in a
non-parametric error memory and selectively invokes a frozen corrective prompt
when similar residual patterns indicate positive utility. ERRORFlow updates no
model weights and explicitly measures the quality--cost trade-off. Using
Qwen3-4B-Instruct-2507, conservative grouped evaluation on AVeriTeC shows that
memory-conditioned routing improves accuracy by 1.50 points and macro-F1 by
1.17 points while routing 11.18\% of examples. On the locked four-label
CLIMATE-FEVER evaluation, always-on correction improves accuracy and macro-F1
by 3.06 and 5.85 points, while selective ERRORFlow obtains smaller significant
gains at 19.93\% routing. On SciFact with open-corpus BM25 evidence, ERRORFlow
routes 18.67\%, improving macro-F1 by 2.58 points with a positive paired
bootstrap interval, although its accuracy interval crosses zero. Retrieval
stratification shows that selection avoids regressions caused by always-on
correction when relevant evidence is missed. Matched-budget comparisons do not
establish superiority over simple confidence routing across external domains,
and conflicting-evidence classification remains weak.
\keywords{Fact verification \and Training-free inference \and Error memory
\and Selective routing \and Retrieval-augmented verification}
\end{abstract}

\section{Introduction}

Fact verification systems must jointly reason over a claim and imperfect
evidence while deciding among support, refutation, insufficient evidence, and
in some datasets conflicting evidence. A second-pass critique can repair some
first-pass errors, but applying it to every example doubles label-scoring work
and can change correct predictions into incorrect ones. The central question
is therefore not only which corrective prompt to use, but when its expected
benefit justifies its cost.

We study this question without updating model parameters. ERRORFlow first
executes a frozen verifier, then retrieves compatible residual outcomes from a
versioned source memory. A selective router invokes one frozen correction
action only when neighboring source outcomes indicate positive utility. The
method records helpful and harmful changes, route rate, and model forward
passes, enabling matched-budget comparisons against confidence, similarity,
label-conditional, static, and random routers.

Our experiments emphasize claim boundaries and reproducibility. Source-domain
evaluation excludes the complete query fold from memory and includes
near-duplicate grouping sensitivity. External target labels are not used for
routing or model selection. We evaluate an ontology-aligned four-label setting
on CLIMATE-FEVER, an adapter-sensitive oracle-rationale setting on SciFact,
and an open-corpus BM25 SciFact condition that exposes retrieval failures.

The main contributions are:
\begin{itemize}
  \item a training-free residual correction pipeline with immutable prompts,
  versioned outcome memory, deterministic routing, and explicit compute
  accounting;
  \item leakage-audited source evaluation with exact policy replay and
  near-duplicate-group sensitivity;
  \item matched-budget controls that separate corrective-action quality from
  routing quality; and
  \item cross-domain and retrieved-evidence analyses that identify both
  transfer gains and failure modes.
\end{itemize}

\section{Related Work}

\paragraph{Fact verification.}
% TODO: cite FEVER, AVeriTeC, SciFact, and CLIMATE-FEVER primary papers.
Prior work combines evidence retrieval with natural-language inference or
large language model reasoning. Direct leaderboard comparisons are outside our
scope because our evidence conditions and frozen verifier differ across
benchmarks.

\paragraph{Test-time correction and compute.}
% TODO: cite self-refinement, critique, verifier, and test-time compute work.
Repeated inference and critique can improve reasoning, but uniform application
obscures regressions and compute cost. ERRORFlow treats correction as a
selective action evaluated against matched inference budgets.

\paragraph{Non-parametric memory.}
% TODO: cite case-based reasoning, retrieval memory, and error-driven systems.
Our memory stores observed action outcomes rather than model parameters. It is
supervised by labeled source outcomes and is therefore not label-free.

\section{Method}

\subsection{Frozen Base Verifier}

Given claim $c$, evidence set $E$, and label set $\mathcal{Y}$, B1 computes the
mean token log probability of each candidate label under a fixed instruction:
\begin{equation}
s_y = \frac{1}{|y|}\sum_{t=1}^{|y|}\log p(y_t\mid c,E,y_{<t}).
\end{equation}
The predicted label is $\hat{y}=\arg\max_y s_y$. Four model forward passes are
used for four-label datasets. We use Qwen3-4B-Instruct-2507 with deterministic
label scoring and seed 2026.

\subsection{Corrective Action}

A3 is a second immutable instruction, \texttt{conflicting\_rescue\_v1}, that
asks the same model to reconsider evidence coverage, polarity, and potential
conflict. A3 uses the same candidate-label scoring rule as B1. Always-on A3
therefore adds four model forward passes per example.

\subsection{Residual Outcome Memory}

Each source-memory record contains the source claim, B1 prediction and score
profile, evidence-length features, the A3 outcome, and observed utility. No
target label enters a route decision. For source evaluation, records from the
query's complete grouped fold are excluded. For external evaluation, the
complete frozen AVeriTeC train memory is used.

Memory retrieval combines claim token Jaccard similarity, B1 label-score
distribution similarity, and evidence-length similarity. R6 retrieves five
compatible records, requires support from at least two, subtracts a fixed
action-cost penalty of 0.1, and routes only when estimated utility is positive.
With one constant-cost action, this penalty cannot change ordering at a fixed
budget; our experiments therefore validate residual outcome-memory routing,
not an independent heterogeneous-cost mechanism.

\subsection{Controls and Cost}

We compare B1, A3-All, confidence-only, similarity-only, and R6 in every
external evaluation. Source ablations additionally include random, margin,
label-conditional, static-pattern, and outcome-memory routers. Selective
controls receive exactly the realized R6 route count. Average model forward
passes equal $4+4r$, where $r$ is the route rate.

\section{Experimental Protocol}

\subsection{Datasets and Evidence Conditions}

\begin{table}[t]
\caption{Evaluation settings. External target labels are used only for final
metrics.}
\label{tab:settings}
\centering
\small
\begin{tabular}{llll}
\toprule
Dataset & $N$ & Labels & Evidence condition \\
\midrule
AVeriTeC train & 3068 & 4 & Human QA evidence; grouped internal evaluation \\
CLIMATE-FEVER & 1535 & 4 & Five official retrieved/annotated candidates \\
SciFact dev & 300 & 3 & Official gold-rationale oracle \\
SciFact dev & 300 & 3 & Open-corpus BM25 top-5 sentences \\
\bottomrule
\end{tabular}
\end{table}

The CLIMATE-FEVER release is evaluated as one locked external population. Its
DISPUTED label maps directly to CONFLICTING. SciFact has no CONFLICTING label;
the predeclared primary adapter maps CONFLICTING predictions to NOT ENOUGH
INFO. We retain raw four-label sensitivity and projection counts.

The SciFact retriever ranks all 5,183 corpus documents with BM25 over title and
abstract, retains five documents, and then ranks their sentences to retain
five. It does not inspect cited document IDs, rationale indices, or labels.
Among 188 verifiable claims, top-5 any-rationale recall is 64.36\% and complete
rationale recall is 62.77\%.

\subsection{Metrics and Statistical Tests}

We report accuracy, macro-F1, per-class F1, helpful and harmful prediction
changes, route rate, and average forward passes. Paired comparisons use exact
two-sided McNemar tests and 5,000-sample paired bootstrap confidence intervals.
Source matched-router comparisons additionally apply Holm correction.

\section{Results}

\begin{table*}[t]
\caption{Main quality--cost results. AVeriTeC uses the conservative
near-duplicate-grouped internal protocol. SciFact results use the predeclared
three-label projection.}
\label{tab:main}
\centering
\small
\begin{tabular}{llrrrr}
\toprule
Evaluation & System & Accuracy & Macro-F1 & Route rate & Avg. passes \\
\midrule
\multirow{2}{*}{AVeriTeC grouped internal}
 & B1 & 0.7402 & 0.5460 & 0.000 & 4.000 \\
 & R6 & 0.7552 & 0.5577 & 0.112 & 4.447 \\
\midrule
\multirow{3}{*}{CLIMATE-FEVER}
 & B1 & 0.5212 & 0.3808 & 0.000 & 4.000 \\
 & A3-All & 0.5518 & 0.4393 & 1.000 & 8.000 \\
 & R6 & 0.5375 & 0.4063 & 0.199 & 4.797 \\
\midrule
\multirow{3}{*}{SciFact oracle}
 & B1 & 0.6067 & 0.5596 & 0.000 & 4.000 \\
 & A3-All & 0.8167 & 0.7822 & 1.000 & 8.000 \\
 & R6 & 0.6367 & 0.6052 & 0.113 & 4.453 \\
\midrule
\multirow{3}{*}{SciFact BM25}
 & B1 & 0.6100 & 0.5778 & 0.000 & 4.000 \\
 & A3-All & 0.6367 & 0.6157 & 1.000 & 8.000 \\
 & R6 & 0.6267 & 0.6035 & 0.187 & 4.747 \\
\bottomrule
\end{tabular}
\end{table*}

\subsection{Source-Domain Routing}

Under conservative token-Jaccard 0.90 grouping, R6 improves AVeriTeC accuracy
by 1.50 points and macro-F1 by 1.17 points while routing 11.18\%. Both paired
bootstrap intervals are positive and exact McNemar gives $p=8.75\times
10^{-5}$. However, the action and router were selected on AVeriTeC train, so
this result is train-selected rather than an independent generalization
estimate. At the matched budget, R6 does not obtain a Holm-significant
macro-F1 advantage over the strongest label-conditional control.

\subsection{External Four-Label Transfer}

On CLIMATE-FEVER, A3-All improves B1 by 3.06 accuracy points and 5.85 macro-F1
points. R6 routes 19.93\% and improves B1 by 1.63 and 2.55 points; both paired
intervals are positive. R6 is not significantly different from the
confidence-only router at the same route count. Per-class analysis further
shows that CONFLICTING F1 changes only from 0.1111 to 0.1202 under A3-All;
most improvement occurs for NOT ENOUGH INFO and REFUTED.

\subsection{SciFact Ontology Sensitivity}

With oracle rationales, projected R6 improves accuracy by 3.00 points and
macro-F1 by 4.56 points. The raw four-label comparison is null. Moreover, A3
produces 96 CONFLICTING outputs under oracle evidence, and projection converts
90 additional examples to correct predictions. This result is therefore
adapter-mediated and cannot support projection-independent transfer.

\subsection{Retrieved-Evidence Robustness}

With open BM25 evidence, R6 improves B1 macro-F1 by 2.58 points with a positive
bootstrap interval, while its 1.67-point accuracy interval crosses zero.
A3-All and R6 are not significantly different. Only five A3 predictions are
projected, so this condition is not driven by mass ontology conversion.

Stratification reveals the mechanism. On 121 verifiable retrieval hits, B1 is
already strong and R6 has net zero correctness change. On 67 retrieval misses,
A3-All loses six net correct predictions, while R6 preserves B1 accuracy. On
112 NEI claims, A3-All gains 15 and R6 gains five. R6 therefore provides
selective harm avoidance rather than recovery of missing facts.

\section{Limitations}

The experiments use one base-model checkpoint. Source results are
train-selected, and source memory uses labeled outcomes; training-free denotes
the absence of weight updates, not the absence of supervision. SciFact oracle
results depend strongly on label projection. CLIMATE-FEVER uses official
retrieved candidates rather than ERRORFlow retrieval, while the open SciFact
retriever is only a lexical baseline. Several external matched-router
comparisons remain statistically inconclusive. Finally, CONFLICTING-class
performance is weak, and constant action cost prevents an independent test of
heterogeneous cost-aware ranking.

\section{Conclusion}

ERRORFlow shows that a frozen corrective action can be invoked selectively
using non-parametric source residual outcomes. The method reduces second-pass
compute and can improve a frozen verifier, but its strongest defensible effect
is selective harm avoidance under imperfect evidence. External results do not
establish universal superiority over simple confidence routing, and ontology
adaptation must be reported explicitly.

\begin{credits}
\subsubsection{Acknowledgments}
% TODO: add funding and acknowledgments.

\subsubsection{Disclosure of Interests}
The authors have no competing interests to declare that are relevant to the
content of this article.
\end{credits}

% TODO: add verified primary references before submission, then uncomment:
% \bibliographystyle{splncs04}
% \bibliography{references}

\end{document}
